\chapter{2008年重庆卷（文）}

\section{单选题}

\begin{problem}
已知 \(\{a_{n}\}\) 为等差数列，\(a_{2} + a_{8} = 12\)，则 \(a_{5}\) 等于（\quad）

\choices
    {\(4\)}
    {\(5\)}
    {\(6\)}
    {\(7\)}
\end{problem}

\begin{answer}
C
\end{answer}

\begin{solution}
等差数列中，首项到第九项关于第五项对称，因此 \(a_{2}+a_{8}=2a_{5}\). 由已知得 \(2a_{5}=12\)，所以 \(a_{5}=6\)，故选 C.
\end{solution}

\begin{problem}
设 \(x\) 是实数，则 \(x > 0\) 是“\(|x| > 0\)”的（\quad）

\choices
    {充分而不必要条件}
    {必要而不充分条件}
    {充要条件}
    {既不充分也不必要条件}
\end{problem}

\begin{answer}
A
\end{answer}

\begin{solution}
若 \(x>0\)，则 \(|x|=x>0\)，所以 \(x>0\) 是 \(|x|>0\) 的充分条件. 但 \(x=-1\) 时 \(|x|=1>0\)，而 \(x>0\) 不成立，所以它不是必要条件，故选 A.
\end{solution}

\begin{problem}
曲线 \(C\)： \(\begin{cases}
x = \cos\theta - 1,\\
y = \sin\theta + 1,
\end{cases}\)（\(\theta\) 为参数）的普通方程为（\quad）

\choices
    {\((x - 1)^2 + (y + 1)^2 = 1\)}
    {\((x + 1)^{2} + (y + 1)^{2} = 1\)}
    {\((x + 1)^{2} + (y - 1)^{2} = 1\)}
    {\((x - 1)^{2} + (y - 1)^{2} = 1\)}
\end{problem}

\begin{answer}
C
\end{answer}

\begin{solution}
由参数方程得 \(x+1=\cos\theta\)，\(y-1=\sin\theta\). 两式平方相加，得到 \((x+1)^{2}+(y-1)^{2}=1\)，故选 C.
\end{solution}

\begin{problem}
若点 \(P\) 分有向线段 \(\overrightarrow{AB}\) 所成的比为 \(-\frac{1}{3}\)，则点 \(B\) 分有向线段 \(\overrightarrow{PA}\) 所成的比是（\quad）

\choices
    {\(-\frac{3}{2}\)}
    {\(-\frac{1}{2}\)}
    {\(\frac{1}{2}\)}
    {\(3\)}
\end{problem}

\begin{answer}
A
\end{answer}

\begin{solution}
按有向线段的比的定义，\(\dfrac{\overrightarrow{AP}}{\overrightarrow{PB}}=-\dfrac{1}{3}\). 设 \(\overrightarrow{AB}=\bs{u}\)，则 \(\overrightarrow{AP}=-\dfrac{1}{2}\bs{u}\)，从而 \(\overrightarrow{PB}=\overrightarrow{AB}-\overrightarrow{AP}=\dfrac{3}{2}\bs{u}\). 因此
\[
\frac{\overrightarrow{PB}}{\overrightarrow{BA}}=\frac{\frac{3}{2}\bs{u}}{-\bs{u}}=-\frac{3}{2}
\]
故选 A.
\end{solution}

\begin{problem}
某校高三年级有男生 \(500\) 人，女生 \(400\) 人. 为了解该年级学生的健康情况，从男生中任意抽取 \(25\) 人，从女生中任意抽取 \(20\) 人进行调查. 这种抽样方法是（\quad）

\choices
    {简单随机抽样法}
    {抽签法}
    {随机数表法}
    {分层抽样法}
\end{problem}

\begin{answer}
D
\end{answer}

\begin{solution}
总体按性别分成男生、女生两层，分别从男生层抽取 \(25\) 人、从女生层抽取 \(20\) 人，且抽取比例分别为 \(\dfrac{25}{500}=\dfrac{20}{400}=\dfrac{1}{20}\). 这是分层抽样，故选 D.
\end{solution}

\begin{problem}
函数 \( y = 10^{x^{2} - 1} \)（\(0 < x \leqslant 1\)）的反函数是（\quad）

\choices
    {\(y = -\sqrt{1 + \lg x}\)（\(x > \frac{1}{10}\)）}
    {\(y = \sqrt{1 + \lg x}\)（\(x > \frac{1}{10}\)）}
    {\(y = -\sqrt{1 + \lg x}\)（\(\frac{1}{10} < x \leqslant 1\)）}
    {\(y = \sqrt{1 + \lg x}\)（\(\frac{1}{10} < x \leqslant 1\)）}
\end{problem}

\begin{answer}
D
\end{answer}

\begin{solution}
原函数的定义域为 \(0<x\leqslant1\). 因为 \(x^{2}\) 在该区间上单调递增，且 \(10^{x^{2}-1}\) 也随指数单调递增，所以原函数在该区间上单调递增，值域为 \(\dfrac{1}{10}<y\leqslant1\). 交换 \(x\)，\(y\)，有 \(x=10^{y^{2}-1}\)，于是 \(y^{2}=1+\lg x\). 因为原函数的自变量满足 \(y>0\)，故反函数为 \(y=\sqrt{1+\lg x}\)，定义域为 \(\dfrac{1}{10}<x\leqslant1\)，故选 D.
\end{solution}

\begin{problem}
函数 \( f(x) = \frac{\sqrt{x}}{x + 1} \) 的最大值为（\quad）

\choices
    {\( \frac{2}{5} \)}
    {\( \frac{1}{2} \)}
    {\( \frac{\sqrt{2}}{2} \)}
    {\(1\)}
\end{problem}

\begin{answer}
B
\end{answer}

\begin{solution}
令 \(t=\sqrt{x}\geqslant0\)，则 \(f(x)=\dfrac{t}{t^{2}+1}\). 由 \((t-1)^{2}\geqslant0\) 得 \(t^{2}+1\geqslant2t\)，所以 \(f(x)\leqslant\dfrac12\). 当 \(t=1\)，即 \(x=1\) 时取等号，故最大值为 \(\dfrac12\)，选 B.
\end{solution}

\begin{problem}
若双曲线 \(\frac{x^2}{3} - \frac{16y^2}{p^2} = 1\) 的左焦点在抛物线 \(y^2 = 2px\) 的准线上，则 \(p\) 的值为（\quad）

\choices
    {\(2\)}
    {\(3\)}
    {\(4\)}
    {\(4 \sqrt{2}\)}
\end{problem}

\begin{answer}
C
\end{answer}

\begin{solution}
双曲线的半实轴平方为 \(3\)，半虚轴平方为 \(\dfrac{p^{2}}{16}\)，所以其焦距的一半满足 \(c^{2}=3+\dfrac{p^{2}}{16}\). 抛物线 \(y^{2}=2px\) 的准线为 \(x=-\dfrac{p}{2}\)，而双曲线左焦点为 \((-c,0)\). 由左焦点在准线上，得 \(c=\dfrac{p}{2}\). 因而
\[
3+\frac{p^{2}}{16}=\frac{p^{2}}{4}
\]
解得 \(p^{2}=16\). 由 \(c=\dfrac{p}{2}>0\) 可知 \(p>0\)，故 \(p=4\)，选 C.
\end{solution}

\begin{problem}
从编号为 \(1\)，\(2\)，\(\ldots\)，\(10\) 的 \(10\) 个大小相同的球中任取 \(4\) 个，则所取 \(4\) 个球的最大号码是 \(6\) 的概率为（\quad）

\choices
    {\( \frac{1}{84} \)}
    {\( \frac{1}{21} \)}
    {\( \frac{2}{5} \)}
    {\( \frac{3}{5} \)}
\end{problem}

\begin{answer}
B
\end{answer}

\begin{solution}
任取 \(4\) 个球的总取法数为 \(\mathrm C_{10}^{4}\). 最大号码为 \(6\) 时必须取到编号 \(6\)，其余三个球从编号 \(1\) 至 \(5\) 中取，取法数为 \(\mathrm C_{5}^{3}\). 因此所求概率为
\[
\frac{\mathrm C_{5}^{3}}{\mathrm C_{10}^{4}}=\frac{10}{210}=\frac{1}{21}
\]
故选 B.
\end{solution}

\begin{problem}
若 \( \left(x+\frac{1}{2x}\right)^{n} \) 的展开式中前三项的系数成等差数列，则展开式中 \( x^{4} \) 项的系数为（\quad）

\choices
    {\(6\)}
    {\(7\)}
    {\(8\)}
    {\(9\)}
\end{problem}

\begin{answer}
B
\end{answer}

\begin{solution}
前三项的系数分别为 \(1\)、\(\dfrac{n}{2}\)、\(\dfrac{n(n-1)}{8}\). 因为前三项存在，\(n\geqslant2\)，它们成等差数列，所以
\[
2\cdot\frac{n}{2}=1+\frac{n(n-1)}{8}
\]
整理得 \(n^{2}-9n+8=0\)，故 \(n=1\) 或 \(n=8\)，结合 \(n\geqslant2\) 得 \(n=8\). 展开式中 \(x^{4}\) 项对应取 \(\left(\dfrac{1}{2x}\right)^{2}\)，其系数为
\[
\mathrm C_{8}^{2}\left(\frac{1}{2}\right)^{2}=7
\]
故选 B.
\end{solution}

\begin{problem}
如图，模块\(\circled{1}\sim\circled{5}\)均由 \(4\) 个棱长为 \(1\) 的小正方体构成，模块\(\circled{6}\)由 \(15\) 个棱长为 \(1\) 的小正方体构成. 现从模块\(\circled{1}\sim\circled{5}\)中选出三个放到模块\(\circled{6}\)上，使得模块\(\circled{6}\)成为一个棱长为 \(3\) 的大正方体. 则下列选择方案中，能够完成任务的为（\quad）

\bitmapfigure[max width=0.82\linewidth]{img/2008/chongqing_liberal/q11_fig1.png}

\choices
    {模块\(\circled{1}\circled{2}\circled{5}\)}
    {模块\(\circled{1}\circled{3}\circled{5}\)}
    {模块\(\circled{2}\circled{4}\circled{6}\)}
    {模块\(\circled{3}\circled{4}\circled{5}\)}
\end{problem}

\begin{answer}
A
\end{answer}

\begin{solution}
模块 \(\circled{6}\) 已有 \(15\) 个小正方体，而棱长为 \(3\) 的大正方体共有 \(27\) 个小正方体，所以还需补入 \(27-15=12\) 个空缺位置. 三个待放模块各有 \(4\) 个小正方体，因此可行的放置必须使它们全部落在空缺位置中，且彼此不重叠.

按图中模块 \(\circled{6}\) 的方向，将模块 \(\circled{5}\) 旋转后放入右下角的空缺处. 此时，最上层的空缺分成左侧的斜线区域和右上方的区域：模块 \(\circled{1}\)（三连方块末端上凸）恰好填入左侧区域，模块 \(\circled{2}\)（\(2\times2\) 方形平板）恰好填入右上方区域. 这三个模块占据的 \(12\) 个位置互不重叠，并且正好补齐全部空缺，所以方案 \(\circled{1}\circled{2}\circled{5}\) 可行.

其余方案中，方案 \(\circled{2}\circled{4}\circled{6}\) 含有模块 \(\circled{6}\)，不符合只能从模块 \(\circled{1}\sim\circled{5}\) 中选取三个模块的条件. 对方案 \(\circled{1}\circled{3}\circled{5}\) 和方案 \(\circled{3}\circled{4}\circled{5}\) 按图逐层核对模块的全部旋转放置位置时，模块 \(\circled{3}\) 的凸出小正方体都会与已占位置重合，或使剩余空缺至少有一格无法填入；因此两方案都不能补齐 \(12\) 个空缺. 故能够完成任务的方案为模块 \(\circled{1}\)、\(\circled{2}\)、\(\circled{5}\)，选 A.
\end{solution}

\begin{problem}
函数 \( f(x) = \frac{\sin x}{\sqrt{5 + 4\cos x}} \)（\(0 \leqslant x \leqslant 2\pi\)）的值域是（\quad）

\choices
    {\( \left[-\frac{1}{4},\frac{1}{4}\right] \)}
    {\( \left[-\frac{1}{3},\frac{1}{3}\right] \)}
    {\( \left[-\frac{1}{2},\frac{1}{2}\right] \)}
    {\(\left[-\frac{2}{3},\frac{2}{3}\right]\)}
\end{problem}

\begin{answer}
C
\end{answer}

\begin{solution}
因为 \(5+4\cos x\geqslant1>0\)，所以函数有意义. 令 \(t=\cos x\in[-1,1]\)，则
\[
f^{2}(x)=\frac{1-t^{2}}{5+4t}=:h(t)
\]
求导得
\[
h'(t)=-\frac{2(2t+1)(t+2)}{(5+4t)^{2}}
\]
在 \([-1,1]\) 上，\(t+2>0\) 且分母恒为正，所以 \(h'(t)>0\) 当 \(-1\leqslant t<-\dfrac12\)，\(h'(t)<0\) 当 \(-\dfrac12<t\leqslant1\). 因此 \(h(t)\) 在 \(t=-\dfrac12\) 处取得最大值，且
\[
h\left(-\frac12\right)=\frac{1-\frac14}{5-2}=\frac14
\]
因此 \(|f(x)|\leqslant\dfrac12\). 当 \(\cos x=-\dfrac12\) 时，取 \(x=\dfrac{2\pi}{3}\) 和 \(x=\dfrac{4\pi}{3}\)，分别得到 \(f(x)=\dfrac12\) 和 \(f(x)=-\dfrac12\). 又 \(f(x)\) 在 \([0,2\pi]\) 上连续，由介值定理可知它取得两端点之间的所有值，故值域为 \(\left[-\dfrac12,\dfrac12\right]\)，选 C.
\end{solution}

\section{填空题}

\begin{problem}
设集合 \(U = \{1, 2, 3, 4, 5\}\)，\(A = \{2, 3, 4\}\)，\(B = \{4, 5\}\)，则 \(A \cap (\complement_U B) =\)\fillinblank{}.
\end{problem}

\begin{answer}
\(\{2,3\}\)
\end{answer}

\begin{solution}
由全集 \(U\) 得 \(\complement_{U}B=\{1,2,3\}\). 因而
\[
A\cap\left(\complement_{U}B\right)=\{2,3,4\}\cap\{1,2,3\}=\{2,3\}
\]
\end{solution}

\begin{problem}
若 \(x > 0\)，则 \(\left(2x^{\frac{1}{4}} + 3^{\frac{3}{2}}\right)\left(2x^{\frac{1}{4}} - 3^{\frac{3}{2}}\right) - 4x^{-\frac{1}{2}}\left(x - x^{\frac{1}{2}}\right) =\)\fillinblank{}.
\end{problem}

\begin{answer}
\(-23\)
\end{answer}

\begin{solution}
利用平方差公式，并注意 \(x>0\)，原式为
\[
\left(2x^{\frac14}\right)^{2}-\left(3^{\frac32}\right)^{2}-4x^{-\frac12}\left(x-x^{\frac12}\right)
=4x^{\frac12}-27-4\left(x^{\frac12}-1\right)=-23
\]
\end{solution}

\begin{problem}
已知圆 \(C: x^{2} + y^{2} + 2x + ay - 3 = 0\)（\(a\) 为实数）上任意一点关于直线 \(l: x - y + 2 = 0\) 的对称点都在圆 \(C\) 上，则 \(a = \)\fillinblank{}.
\end{problem}

\begin{answer}
\(-2\)
\end{answer}

\begin{solution}
将圆方程配方，得
\[
\left(x+1\right)^{2}+\left(y+\frac{a}{2}\right)^{2}=4+\frac{a^{2}}{4}
\]
所以圆心为 \(O\left(-1,-\dfrac{a}{2}\right)\). 一个圆关于直线反射后仍与自身重合，当且仅当圆心在该直线上. 代入直线 \(x-y+2=0\)，得
\[
-1+\frac{a}{2}+2=0
\]
即 \(a=-2\).
\end{solution}

\begin{problem}
某人有 \(3\) 种颜色的灯泡（每种颜色的灯泡足够多），要在如图所示的 \(6\) 个点
\(A\)、\(B\)、\(C\)、\(A_1\)、\(B_1\)、\(C_1\) 上各装一个灯泡，要求同一条线段两端的灯泡不同色，

\bitmapfigure[max width=0.45\linewidth]{img/2008/chongqing_liberal/q16_fig1.png}

则不同的安装方法共有\fillinblank{} 种. （用数字作答）
\end{problem}

\begin{answer}
12
\end{answer}

\begin{solution}
图中的线段构成一个三棱柱的棱图：上、下两个三角形的三个顶点各两两相连，并有三条对应的竖棱相连. 先给上三角形 \(ABC\) 安装灯泡，三个顶点两两相连，因此三种颜色必须各用一次，有 \(3!=6\) 种方法.

固定上三角形的一种颜色后，下三角形 \(A_{1}B_{1}C_{1}\) 也必须三色各异；同时 \(A_{1}\)、\(B_{1}\)、\(C_{1}\) 不能分别与 \(A\)、\(B\)、\(C\) 同色. 在三种颜色的排列中，满足这三个限制的只有两个循环排列. 所以总安装方法数为 \(6\times2=12\) 种.
\end{solution}

\section{解答题}

\begin{problem}
设 \(\triangle ABC\) 的内角 \(A\)，\(B\)，\(C\) 的对边分别为 \(a\)，\(b\)，\(c\). 已知 \(b^2 + c^2 = a^2 + \sqrt{3}bc\)，求：
\begin{enumerate}
\item \(A\) 的大小；
\item \(2\sin B\cos C - \sin (B - C)\) 的值.
\end{enumerate}
\end{problem}

\begin{solution}
\begin{enumerate}
\item 由余弦定理，\(a^{2}=b^{2}+c^{2}-2bc\cos A\). 与已知等式比较，得 \(2bc\cos A=\sqrt{3}bc\). 因为 \(b\)，\(c>0\)，所以 \(\cos A=\dfrac{\sqrt{3}}{2}\). 又 \(0<A<\pi\)，故 \(A=\dfrac{\pi}{6}\).
\item 利用两角差公式，
\[
2\sin B\cos C-\sin\left(B-C\right)=\sin B\cos C+\cos B\sin C=\sin\left(B+C\right)
\]
由 \(A+B+C=\pi\) 及 \(A=\dfrac{\pi}{6}\)，得 \(B+C=\dfrac{5\pi}{6}\)，所以原式等于 \(\sin\dfrac{5\pi}{6}=\dfrac12\).
\end{enumerate}
\end{solution}

\begin{problem}
在每道单项选择题给出的 \(4\) 个备选答案中，只有一个是正确的. 若对 \(4\) 道选择题中的每一道都任意选定一个答案，求这 \(4\) 道题中：
\begin{enumerate}
\item 恰有两道题答对的概率；
\item 至少答对一道题的概率.
\end{enumerate}
\end{problem}

\begin{solution}
\begin{enumerate}
\item 各道题的选择相互独立，每道题答对的概率为 \(\dfrac14\)，答错的概率为 \(\dfrac34\). 四道题中恰有两道答对的概率为
\[
\mathrm C_{4}^{2}\left(\frac14\right)^{2}\left(\frac34\right)^{2}=\frac{27}{128}
\]
\item “至少答对一道题”是“没有答对任何题”的对立事件. 所以所求概率为
\[
1-\left(\frac34\right)^{4}=1-\frac{81}{256}=\frac{175}{256}
\]
\end{enumerate}
\end{solution}

\begin{problem}
设函数 \( f(x) = x^{3} + ax^{2} - 9x - 1 \)（\(a < 0\)）. 若曲线 \( y = f(x) \) 的斜率最小的切线与直线 \( 12x + y = 6 \) 平行，求：
\begin{enumerate}
\item \(a\) 的值；
\item 函数 \( f(x) \) 的单调区间.
\end{enumerate}
\end{problem}

\begin{solution}
\begin{enumerate}
\item 曲线在横坐标 \(x\) 处的切线斜率为 \(f'(x)=3x^{2}+2ax-9\). 这是关于 \(x\) 的开口向上的二次函数，其最小值为 \(-\dfrac{a^{2}}{3}-9\). 直线 \(12x+y=6\) 的斜率为 \(-12\)，所以
\[
-\frac{a^{2}}{3}-9=-12
\]
得 \(a^{2}=9\). 结合 \(a<0\)，得 \(a=-3\).
\item 此时 \(f'(x)=3x^{2}-6x-9=3(x+1)(x-3)\). 因此 \(f'(x)>0\) 时，\(x\in(-\infty,-1)\cup(3,+\infty)\)，函数在这两个区间上递增；\(f'(x)<0\) 时，\(x\in(-1,3)\)，函数在该区间上递减.
\end{enumerate}
\end{solution}

\begin{problem}
如图，\(\alpha\) 和 \(\beta\) 为平面，\(\alpha \cap \beta = l\)，\(A \in \alpha\)，\(B \in \beta\)，\(AB = 5\)，\(A\)，\(B\) 在棱 \(l\) 上的射影分别为 \(A'\)，\(B'\)，\(AA' = 3\)，\(BB' = 2\). 若二面角 \(\alpha - l - \beta\) 的大小为 \(\frac{2\pi}{3}\)，求：

\begin{enumerate}
\item 点 \(B\) 到平面 \(\alpha\) 的距离；
\item 异面直线 \(l\) 与 \(AB\) 所成的角. （用反三角函数表示）
\end{enumerate}

\bitmapfigure[max width=0.54\linewidth]{img/2008/chongqing_liberal/q20_fig1.png}
\FloatBarrier
\end{problem}

\begin{solution}
\begin{enumerate}
\item 在垂直于 \(l\) 的截面中，设从 \(l\) 分别指向 \(A\)、\(B\) 的横截面向量为 \(\bs{u}\)、\(\bs{v}\)，则 \(|\bs{u}|=3\)，\(|\bs{v}|=2\)，且它们的夹角为二面角 \(\dfrac{2\pi}{3}\). 向量 \(\overrightarrow{AB}\) 垂直于 \(l\) 的分量为 \(\bs{v}-\bs{u}\)，故
\[
|\bs{v}-\bs{u}|^{2}=3^{2}+2^{2}-2\cdot3\cdot2\cos\frac{2\pi}{3}=19
\]
在垂直于 \(l\) 的截面内，\(AA'\) 与 \(BB'\) 的夹角为 \(\dfrac{2\pi}{3}\)，故直线 \(BB'\) 与平面 \(\alpha\) 所成的锐角为 \(\dfrac{\pi}{3}\)，所以点 \(B\) 到平面 \(\alpha\) 的距离为
\[
2\sin\frac{\pi}{3}=\sqrt{3}
\]
\item 由 \(AB=5\)，沿 \(l\) 的分量长度为 \(A'B'\)，有 \(AB^{2}=A'B'^{2}+19\)，从而 \(A'B'^{2}=6\). 设异面直线 \(l\) 与 \(AB\) 所成的锐角为 \(\theta\)，则 \(AB\) 垂直于 \(l\) 的分量长度为 \(5\sin\theta\)，所以
\[
\sin\theta=\frac{\sqrt{19}}{5}
\]
故 \(\theta=\arcsin\dfrac{\sqrt{19}}{5}\).
\end{enumerate}
\end{solution}

\begin{problem}
如图，\(M(-2,0)\) 和 \(N(2,0)\) 是平面上的两点，动点 \(P\) 满足： \(\left||PM|-|PN|\right| = 2\).

\begin{enumerate}
\item 求点 \(P\) 的轨迹方程；
\item 设 \(d\) 为点 \(P\) 到直线 \(l: x = \frac{1}{2}\) 的距离，若 \(|PM| = 2|PN|^2\)，求 \(\frac{|PM|}{d}\) 的值.
\end{enumerate}

\bitmapfigure[max width=0.46\linewidth]{img/2008/chongqing_liberal/q21_fig1.png}
\FloatBarrier
\end{problem}

\begin{solution}
\begin{enumerate}
\item 设 \(|PM|=m\)，\(|PN|=n\). 条件为 \(|m-n|=2\)，而两定点 \(M\)，\(N\) 的距离为 \(4\)，所以轨迹是焦点为 \(M\)，\(N\)、实轴长为 \(2\) 的双曲线. 其半实轴 \(a=1\)，焦半距 \(c=2\)，故 \(b^{2}=c^{2}-a^{2}=3\)，轨迹方程为
\[
\frac{x^{2}}{1^{2}}-\frac{y^{2}}{3}=1
\]
\item 由 \(|PM|=2|PN|^{2}\) 及 \(|m-n|=2\) 分情况讨论. 若 \(m-n=2\)，则 \(2n^{2}-n-2=0\)，由 \(n>0\) 得 \(n=\dfrac{1+\sqrt{17}}{4}\). 若 \(n-m=2\)，则 \(2n^{2}-n+2=0\)，无实根，所以只能有 \(m-n=2\).

由两点距离平方之差，\(m^{2}-n^{2}=8x\)，故 \(x=\dfrac{(m-n)(m+n)}{8}=\dfrac{m+n}{4}=\dfrac{n+1}{2}\). 因而到直线 \(x=\dfrac12\) 的距离为 \(d=\dfrac{n}{2}\). 于是
\[
\frac{|PM|}{d}=\frac{m}{d}=\frac{n+2}{n/2}=2+\frac4n=1+\sqrt{17}
\]
\end{enumerate}
\end{solution}

\begin{problem}
设各项均为正数的数列 \(\{a_{n}\}\) 满足 \(a_1 = 2\)，\(a_n = a_{n+1}^{\frac{3}{2}} a_{n+2}\)（\(n \in \N^*\)）.
\begin{enumerate}
\item 若 \(a_2 = \frac{1}{4}\)，求 \(a_3\)，\(a_4\)，并猜想 \(a_{2008}\) 的值 （不需证明）；
\item 若 \(2\sqrt{2} \leqslant a_1a_2 \cdots a_n < 4\) 对 \(n \geqslant 2\) 恒成立，求 \(a_2\) 的值.
\end{enumerate}
\end{problem}

\begin{solution}
\begin{enumerate}
\item 由递推式在 \(n=1\) 时得
\[
2=\left(\frac14\right)^{\frac32}a_{3}=\frac18a_{3}
\]
所以 \(a_{3}=16\). 再令 \(n=2\)，得
\[
\frac14=16^{\frac32}a_{4}=64a_{4}
\]
所以 \(a_{4}=\dfrac{1}{256}\). 由此得到前几项的指数为 \(1\)，\(-2\)，\(4\)，\(-8\)，故猜想 \(a_{n}=2^{(-2)^{n-1}}\)，从而 \(a_{2008}=2^{(-2)^{2007}}\).
\item 令 \(c_{n}=\log_{2}a_{n}\). 因为各项为正，递推式等价于
\[
c_{n}=\frac32c_{n+1}+c_{n+2}
\]
其特征方程为 \(r^{2}+\dfrac32r-1=0\)，根为 \(\dfrac12\)、\(-2\)，所以
\[
c_{n}=A\left(\frac12\right)^{n-1}+B(-2)^{n-1}
\]
设 \(S_{n}=c_{1}+c_{2}+\cdots+c_{n}\)，则
\[
S_{n}=2A\left(1-2^{-n}\right)+\frac{B}{3}\left(1-(-2)^{n}\right)
\]
题设给出 \(\dfrac32\leqslant S_{n}<2\)（\(n\geqslant2\)），所以 \(S_n\) 有界，必须有 \(B=0\). 又 \(c_{1}=\log_{2}2=1\)，故 \(A=1\)，从而 \(c_{2}=\dfrac12\). 因此 \(a_{2}=2^{c_{2}}=\sqrt{2}\). 反之取 \(a_{2}=\sqrt{2}\) 时，\(S_{n}=2(1-2^{-n})\)，对所有 \(n\geqslant2\) 都满足 \(\dfrac32\leqslant S_{n}<2\)，故该值确实符合题设条件.
\end{enumerate}
\end{solution}
